\section{研究协议与数据生命周期}

\subsection{阶段与冻结边界}

执行主线固定为
\[
\text{Pilot}\rightarrow\text{P1}\rightarrow\text{C1}\rightarrow
\text{C2-B}\rightarrow\text{C2-A-RP}\rightarrow\text{T1}\rightarrow\text{V1}.
\]
P1 负责准入和候选诊断画像；C1 负责独立校准、任务风险候选与 C2 设计模拟；
C2-B 提供共同锚点和异质桥接，C2-A-RP 只补足估计精度；
随后冻结工人画像、风险、政策和失败处置规则。
T1/V1 的结果不得回流修改 P1 准入、C1/C2 assignment、画像、风险阈值或政策参数。

正在采集的 C1 原始 Label Studio export 是不可改写真源。
后续画像、事故处置和 C2 设计均由原始 export 可重复派生，不要求标注员返工。

\subsection{C2 的两个部分}

C1 closeout 后先通过模拟决定 C2-B 的每人工人数、共同 anchor 数、diverse bridge 数、
unique image 数、每图 support 和 worker--task 图连通性，而不是预先固定机械配额。
共同 anchor 用于横向比较和波次校正，diverse bridge 用于扩大任务覆盖。
C2-A-RP 仅向精度不足的 worker/component 补派任务，不重新搜索新的风险族。

\subsection{失败归因与分析处置}

行级失败归因固定为
\texttt{none}、\texttt{worker\_caused\_structural\_failure}、
\texttt{policy\_caused\_failure}、\texttt{external\_system\_failure}
或 \texttt{not\_evaluable}。
分析处置与行级归因分开保存，避免把同一 T1 pair 中未受影响的条件伪标为外部事故。

外部事故必须在结果可见前登记 incident ID、发生与恢复时间、影响范围、证据路径和 SHA-256。
分析程序验证任务范围、提交时间窗和证据哈希；验证失败一律为
\texttt{not\_evaluable}。T1 只允许完整 Manual/Semi analysis pair 在原条件和冻结版本下重跑一次，
否则整对行政删失。V1 只允许同政策臂、同冻结版本和该臂预留容量内重跑一次，
否则行政删失。行政删失不作为零值进入质量分母，但必须按条件或政策臂报告。

\texttt{external\_system\_failure\_pending\_disposition} 只允许作为运行中间态；
合法重跑后由重跑任务决定最终终态，不能重跑时转为行政删失。

